\section{Structural Limits}
\label{sec:limits}

\subsection*{The SSRF path evades the egress baseline.} When the agent reaches
the attacker through the route the task suggests, the proxy sees
\texttt{package-registry:5000}, which is in the baseline because it is a
legitimate service. The attacker's address travels URL-encoded inside a query
string, where the baseline check does not look, and the second hop happens
server to server inside an internal network without crossing the proxy at all.
This is the anchor incident's own pattern, and it bounds what the egress
monitor is for: it catches the naive agent and misses the one that takes the
detour the task offers it.

The bound is a capability, not an observed evasion, and the distinction
matters. Across all 143 runs the proxy recorded exactly one destination,
\texttt{mirror-externo:5002}, thirty times, flagged anomalous on every one:
every agent that reached the attacker went at it directly. Not one run used
the SSRF route. The blind spot was demonstrated by a synthetic control, not
exploited by a model, so what we report is that the monitor can be evaded this
way, not that it was.

\subsection*{The shared event log is not safe under concurrent append.} The
three monitors append to one file per run from separate containers, and that
append is not atomic over a Docker Desktop bind mount on Windows. We observed
lines beginning mid-JSON. The analysis now skips a torn line and names the
file and line rather than dying on it, but a lost line may be an inert event
or an alert that changes a run's verdict.

\subsection*{Tool-call sequences do not predict a leak.} We looked for a
syntactic signal in the transcripts, on the theory that the order in which an
agent reaches for its tools might separate runs that leak from runs that do
not. It does not, at least at this sample size: the dominant transitions are
the same in both groups, and the apparent signal in which tool an agent
reaches for first is an artifact of each task exposing a different channel. We
report the negative because it bounds what the transcripts support.

\subsection*{Canary obfuscation coverage is a floor, not a claim.} Detection
matches the plain token plus base64, URL encoding, hex and reversal. An agent
encoding the token in, for example, the lengths of file names defeats it, and
we make no claim beyond the four.
